\section{The grid structures}\label{sec:structures}

The instances used in the lower bound attach finite groups to the edges
of a three-dimensional box.  At each vertex we retain the tuples of edge
states satisfying a prescribed divergence equation.  The group law will
supply the nonabelian automorphisms used in the support argument; the
query just defined will detect an obstruction to the divergence equations.

\subsection{Grid orientations and edge groups}

For an integer $L\ge2$, let $G=(V,E)$ be the box grid with
\[
 V=\{0,\ldots,L-1\}^3.
\]
Each edge is oriented in its positive coordinate direction.  Put
$\epsilon_{ve}=1$ if $v$ is the head of $e$, $\epsilon_{ve}=-1$ if it is
the tail, and $\epsilon_{ve}=0$ otherwise.  The incidence map is
\[
 \partial:\mathbb F^E\longrightarrow\mathbb F^V,
 \qquad (\partial z)_v=\sum_{e\in E}\epsilon_{ve}z_e.
\]
Let $F$ be the set of unit square faces.  A face parallel to coordinate
directions $i<j$ is traversed first in direction $+i$, then $+j$, then
$-i$, then $-j$.  Its edge vector $c^f\in\mathbb F^E$ has coefficient
$1$ when this traversal follows an edge orientation, $-1$ when it opposes
that orientation, and zero off the face.  At each vertex of the face one
traversed edge enters and one leaves, so
\begin{equation}\label{eq:face-cycle}
 \partial c^f=0.
\end{equation}
Write $F(e)=\{f\in F:c^f_e\ne0\}$, and let $F(v)$ be the set of faces
containing $v$.  We shall use
\[
 |F(e)|\le4,\qquad f(v):=|F(v)|\le12,
 \qquad d(v):=|\{e:e\ni v\}|\le6.
\]
In particular, $d(v)\ge1$.

For vectors in $\mathbb F^2$, define the alternating form
\[
 \omega((a,b),(a',b'))=ab'-ba'.
\]
The edge groups use the standard finite Heisenberg multiplication
\cite[Section~2.1]{gurevichHadani2009}, with the face-boundary signs
specifying its alternating form.
The group $P_e$ attached to an edge $e$ has underlying set
$(\mathbb F^2)^{F(e)}\times\mathbb F$.  Its elements are written $(u,z)$,
where $u=(u_f)_{f\in F(e)}$, and its product is
\begin{equation}\label{eq:edge-product}
 (u,z)(u',z')=
 \left(u+u',\ z+z'+\frac12
          \sum_{f\in F(e)}c^f_e\omega(u_f,u'_f)\right).
\end{equation}
Here $1/2=2$ in $\mathbb F_3$.  To verify the group law, denote the
bilinear correction term by $B_e(u,u')$.  Bilinearity gives
\[
 B_e(u,u')+B_e(u+u',u'')
   =B_e(u',u'')+B_e(u,u'+u''),
\]
which is exactly associativity.  Alternation gives $B_e(u,u)=0$; hence
$(0,0)$ is the identity and $(u,z)^{-1}=(-u,-z)$.
In particular, the last coordinate of $(u,z)^{-1}(u',z')$ is
\begin{equation}\label{eq:edge-difference}
 z'-z-\frac12\sum_{f\in F(e)}c^f_e\omega(u_f,u'_f).
\end{equation}

\subsection{Atoms and binary relations}

Fix $b=(b_v)_{v\in V}\in\mathbb F^V$.  The universe of $A_b$ is the
disjoint union of two kinds of atoms.  There is an \emph{edge atom}
$(e,u,z)$ for each edge $e$ and each $(u,z)\in P_e$.  At each vertex $v$
there is a \emph{configuration atom} for every tuple
\[
 a=\bigl(v,((u_e,z_e))_{e\ni v}\bigr)
\]
satisfying
\begin{equation}\label{eq:configuration}
 \begin{aligned}
 u_{e,f}&=u_{e',f}
       &&\text{for the two edges $e,e'$ of each face $f\in F(v)$ at $v$},\\
 \sum_{e\ni v}\epsilon_{ve}z_e&=b_v.
 \end{aligned}
\end{equation}
The coordinates in these descriptions label atoms; they are not members
of the atoms as sets.  The two kinds of atoms are tagged so they are
disjoint.

Figure~\ref{fig:unit-face} illustrates the signed face boundary in
\eqref{eq:face-cycle} and the shared-coordinate condition in
\eqref{eq:configuration}.
\begin{figure}[t]
\centering
\begin{tikzpicture}[x=1cm,y=1cm,>=Stealth,font=\small]
  \draw[->] (0,0) -- (2.2,0);
  \draw[->] (2.2,0) -- (2.2,2.2);
  \draw[->] (0,2.2) -- (2.2,2.2);
  \draw[->] (0,0) -- (0,2.2);
  \foreach \x/\y in {0/0,2.2/0,2.2/2.2,0/2.2}
    \fill (\x,\y) circle (1.2pt);
  \node[below left] at (0,0) {$v$};
  \node at (1.1,1.1) {$f$};
  \node[above] at (1.1,0) {$+1$};
  \node[left] at (2.2,1.1) {$+1$};
  \node[below] at (1.1,2.2) {$-1$};
  \node[right] at (0,1.1) {$-1$};
  \node[above] at (1.1,2.2) {$+i$};
  \node[right] at (2.2,1.1) {$+j$};
  \node[below] at (1.1,-0.08) {$e,\ u_{e,f}$};
  \node[left] at (-0.08,1.1) {$e',\ u_{e',f}$};
  \node at (1.1,-0.72) {$u_{e,f}=u_{e',f}$};
\end{tikzpicture}

\caption{A geometric grid face $f$ in directions $i<j$. Arrows follow the
global positive-coordinate orientations; the labels $+1,+1,-1,-1$ are
the coefficients of $c^f$, so their signed divergence is zero at each
corner. At the marked vertex $v$, the two incident states of one
configuration satisfy $u_{e,f}=u_{e',f}$. The drawing shows the geometric
cell underlying these coordinates; the atoms are the states and
configurations defined in the text.}
\label{fig:unit-face}
\end{figure}

\begin{lemma}\label{lem:universe-size}
For each $L\ge2$, the size of $A_b$ is independent of $b$ and equals
\begin{equation}\label{eq:universe-size}
 N_L=\sum_{e\in E}3^{2|F(e)|+1}
       +\sum_{v\in V}3^{2f(v)+d(v)-1}.
\end{equation}
In particular,
\[
 L^3\le N_L\le (3^{10}+3^{29})L^3.
\]
Every vertex has at least one configuration atom.
\end{lemma}
\begin{proof}
There are $3^{2|F(e)|+1}$ states on $e$.  At a vertex $v$, each face
$f\in F(v)$ meets precisely two incident edges.  Their face coordinates
are identified by \eqref{eq:configuration}, independently of all other
faces.  The choices of $u$ therefore contribute $3^{2f(v)}$ possibilities.
The divergence condition is a single linear equation in $d(v)$ variables
with a nonzero coefficient, so it has $3^{d(v)-1}$ solutions for every
$b_v$.  This proves \eqref{eq:universe-size} and positivity.  There are
$L^3$ vertices and $3L^2(L-1)\le3L^3$ edges.  The displayed bounds follow
from $|F(e)|\le4$, $f(v)\le12$, and $d(v)\le6$.
\end{proof}

The vocabulary $\tau$ is interpreted as follows; each relation is false
on all pairs not specified below.
\begin{itemize}
\item $\mathsf{Ed}$ and $\mathsf{Cf}$ are the diagonals of the edge atoms
      and configuration atoms, respectively.
\item $\mathsf{EB}$ relates edge atoms on the same edge.  The relation
      $\mathsf{VB}$ relates configuration atoms at the same vertex.
\item $\mathsf I(a,x)$ holds when $a$ is a configuration at a vertex $v$
      and $x=(e,u_e,z_e)$ is the state selected by $a$ on an incident
      edge $e$.
\item For $\delta\in\mathbb F$, the relation
      $\mathsf Z_\delta((e,u,z),(e,u',z'))$ holds precisely when
      \eqref{eq:edge-difference} equals $\delta$.
\end{itemize}
Thus $\mathsf I$ goes from configurations to edge atoms; it is not made
symmetric.  The coordinate labels used to construct $A_b$ are not extra
symbols in its input vocabulary.  In particular, no ordering of the grid
or of its blocks is supplied to a program.

\subsection{Opposite answers}

Fix any vertex $v_*\in V$, and put
\[
 b^0=0,\qquad b^1=\mathbf1_{v_*}.
\]
These two charge vectors give structures of the same size by
Lemma~\ref{lem:universe-size}.

\begin{proposition}\label{prop:opposite-answers}
For every $L\ge2$, the structure $A_{b^0}$ satisfies $Q$ and
$A_{b^1}$ does not.  More generally, $Q(A_b)$ implies
$\sum_{v\in V}b_v=0$ in $\mathbb F$.
\end{proposition}
\begin{proof}
For clarity, let $X(v)$ denote the configuration block at a grid vertex
$v$.  For any $t\in X(v)$, the set $X_t$ in the definition of the query is
exactly $X(v)$.  For an edge atom $y$ on $e$, the test
\eqref{eq:query-incidence-test} holds precisely when $e\ni v$: every
configuration at $v$ selects one atom on each incident edge, and
$\mathsf{EB}$ identifies exactly the edge blocks.  If $e\ni v$, the
coefficient $C(y,a)$ has exactly one contributing selected edge atom,
and exactly one of its three $\mathsf Z_\delta$ relations holds.

On $A_{b^0}$, each vertex has a configuration with all coordinates zero.
Give it weight $\lambda_a=1$ and give all other configurations weight zero.
For $y=(e,u,z)$ set $\mu_y=-z$.  The normalization equations hold.
The selected neighbor of every nonzero-weight configuration on $e$ is
$(e,0,0)$, whose group difference from $y$ has last coordinate $-z$ by
\eqref{eq:edge-difference}.  Thus all equations
\eqref{eq:query-consistency} also hold.

Now suppose the query system on $A_b$ has any solution.  The weights
$\lambda_a$ are arbitrary field elements; we do not assume that they
select a configuration.  For each edge choose, only for this proof, its
zero atom $y_e=(e,0,0)$.  By \eqref{eq:edge-difference}, for $e\ni v$ and
$a\in X(v)$ we have $C(y_e,a)=z_e(a)$.  The equations of the query give
\[
 \mu_{y_e}=\sum_{a\in X(v)}\lambda_a z_e(a),\qquad
 \sum_{a\in X(v)}\lambda_a=1.
\]
Multiply the first identity by $\epsilon_{ve}$, sum over the edges at
$v$, and use \eqref{eq:configuration}:
\[
 \sum_{e\ni v}\epsilon_{ve}\mu_{y_e}
 =\sum_{a\in X(v)}\lambda_a
       \sum_{e\ni v}\epsilon_{ve}z_e(a)
 =b_v\sum_{a\in X(v)}\lambda_a=b_v.
\]
Summing this identity over all vertices cancels each edge contribution,
because an oriented edge has one head and one tail.  Hence
$\sum_v b_v=0$.  For $b=b^1$ that sum is $1$, which is impossible.
\end{proof}

The remainder of the argument shows that a fixed choiceless computation
cannot distinguish these two structures when $L$ is sufficiently large.
The equality of their universe sizes and the fixed binary vocabulary
will be used throughout that argument.
